\documentclass{article}
\usepackage[left=2cm, right=2cm, top=2cm, bottom=2cm]{geometry}
\usepackage{graphicx}
\usepackage{amsmath}
\usepackage{amssymb}
\usepackage{amsthm}
\usepackage{fancyhdr}
\usepackage{verbatim}
\usepackage{listings}
\usepackage{xcolor}
\usepackage{pdfpages}
\usepackage{cancel}
\usepackage{physics}
\usepackage{esint}
\usepackage{braket}

\lstset{
    language=C++,
    basicstyle=\ttfamily\footnotesize,
    keywordstyle=\color{blue},
    commentstyle=\color{green},
    stringstyle=\color{red},
    numbers=left,
    numberstyle=\tiny\color{gray},
    frame=single,
    breaklines=true,
    captionpos=b,
}


\title{PHYS 427: Lecture \# 6}
\author{Cliff Sun}

\newtheorem{theorem}{Theorem}[section]
\newtheorem{lemma}[theorem]{Lemma}
\newtheorem{definition}[theorem]{Definition}
\newtheorem{conjecture}[theorem]{Conjecture}
\newtheorem{proposition}[theorem]{Proposition}
\newtheorem{corollary}[theorem]{Corollary}
\newtheorem{one minute paper}[theorem]{One Minute Paper}

\pagestyle{fancy}
\lhead{\textbf{\thepage}\ \ \nouppercase{\rightmark}}
\chead{PHYS 427: Lecture \# 6}
\rhead{Cliff Sun}

\begin{document}

\maketitle

\section*{Lecture Span}
\begin{enumerate}
    \item Equipartition Theorem
\end{enumerate}

\section*{Recap:}

We got to a one particle in a box (low and high temp limits). In the high temp limit, replace sum with integral. We get a Gaussian integral for $Z$. As 

\begin{equation}
    Z(T) \approx \int_{0}^{\infty}dn_x \exp(- \alpha_x n_x^2)\int_{0}^{\infty}dn_y \exp(- \alpha_y n_y^2)\dots
\end{equation}

Such that 

\begin{equation}
    \alpha_x = \frac{\hbar^2 \pi^2}{2mL_x^2 k_B T}
\end{equation}

The integral is then 

\begin{equation}
    Z_1(T) = \frac{\pi^{3/2}}{2^3 (\alpha_x \alpha_y \alpha_z)^{1/2}}
\end{equation}

You can also plug in $\alpha$, 

\begin{equation}
    Z_1(T) = \left( \frac{m k_B T}{2\pi\hbar^2} \right)^{3/2}V
\end{equation}

For $\alpha_i \ll 1$. Therefore, $k_B T \gg \hbar^2 \pi^2/2mL_x^2$ (same for other directions). 

All above work is in \textbf{the high temperature limit}. 

\begin{equation}
    Z_{tot} = \frac{Z_1^N}{N!}
\end{equation}

Here, use the previously calculated $Z_1$. Calculate $U$ for an ideal gas, then start at the partition function:

Note that $Z_1 \sim \left( \beta^{-3/2} \right)^{N}$

\begin{equation}
    U = -\partial_{\beta}\ln Z_{tot} = -N\partial_{\beta}\ln\left( \beta^{-3/2} + \text{ terms independent of $\beta$} \right) 
\end{equation}

\begin{equation}
    = \frac{3}{2}N\frac{1}{\beta} = \frac{3}{2}Nk_B T = U
\end{equation}

The de-Broglie wavelength $$\lambda = \frac{h}{p} \iff p = h/\lambda$$

Then, 

\begin{equation}
    \langle \epsilon \rangle = \langle p^2/2m \rangle = \frac{3}{2}k_B T 
\end{equation}

Then 

\begin{equation}
    \langle \lambda \rangle = \frac{h}{\sqrt{2m \langle e \rangle}} = \sqrt{\frac{4\pi^2 \hbar^2}{3m k_B T}}
\end{equation}

\subsection*{Quantum Density}

\begin{equation}
    n_Q = \frac{1}{\sqrt{2\pi \hbar^2/m k_B T}^{3/2}}
\end{equation}

For a cube, the edge length

\begin{equation}
    \sqrt{\frac{2\pi \hbar^2}{mk_B T}} \approx \lambda \cdot \mathcal{O}(1)
\end{equation}

\begin{equation}
    n_Q \approx \frac{\text{constant of order 1}}{\langle \lambda \rangle^3}
\end{equation}

\begin{equation}
    Z_1(T) = n_Q(T)V
\end{equation}

Consider Helium at room temperature. That is $T \approx 300$K and atmospheric pressure. Then, 

\begin{equation}
    n_Q(T) \approx 0.8 \times 10^{25}\frac{1}{\text{cm}^3}
\end{equation}

This is about 1 atom per 0.5 Angstrom cubed. The typical density is 

\begin{equation}
    n_{He} = 2.5 \times 10^{19} \frac{1}{\text{cm}^3}= \frac{N}{V} 
\end{equation}

This is about 1 atom per 34 Angstrom cubed. Therefore,

\begin{equation}
    n_{He} \ll n_Q
\end{equation}

This is the \underline{classical regime}.

\section*{Equipartition Theorem}
Considering quadratic degrees of freedom. This means 

\begin{equation}
    \frac{p^2}{2m} = \sum_i \frac{p_i^2}{2m}
\end{equation}

This is 3 degrees of freedom (also quadratic!). Then, 

\begin{equation}
    \langle \epsilon_x \rangle = \langle p_x^2/2m\rangle = \frac{1}{2}k_B T
\end{equation}

Every degree of freedom contributes $(1/2) k_B T$ to $\langle \epsilon \rangle$ in the \textbf{high temperature limit}.

In an harmonic oscillator, there are two degrees of freedom. Since you have $1/2 \cdot k x^2$ and $p_x^2/2m$. For a solid in 3D, then 

\begin{equation}
    \langle \epsilon \rangle = 3Nk_B T
\end{equation}

\subsection*{General statement at high temperatures}

Consider a diatomic gas, CO. Then, we have vibrational (spring system between the two atoms) and translational energy.

The total energy of the molecule is then 

\begin{equation}
    \epsilon_{tot} = \epsilon_{trans} + \epsilon_{rot} + \epsilon_{vib}
\end{equation}

Then,

\begin{equation}
    Z_{tot} = \sum_{\text{all states}}e^{-\beta e_{tot}} = \sum_{e_{trans}}e^{-\beta\epsilon_{trans}} \cdot \sum \dots
\end{equation}

We can then define $Z_{trans}$, $Z_{rot}$, etc.

Consider quantum rigid motor, angular momentum. The total angular momentum $J$ is quantized s.t. its magnitude is 

\begin{equation}
    J = |j| = \hbar\sqrt{j(j+1)}, \quad j = 0,1,2,\dots
\end{equation}
\begin{equation}
    J_z = m_j \hbar, \quad m_j = -j, -j+1, \dots, j-1, j
\end{equation}

\begin{equation}
    \epsilon_{rot} = \frac{|j|^2}{2I} = \frac{j(j+1)\hbar^2}{2I}
\end{equation}

Here, $I$ is the moment of inertia. 

\begin{equation}
    Z_{rot} = \sum_{J, J_z}e^{-\beta \epsilon_{rot}} = \sum_{j=0}^{\infty}(2j + 1)e^{-\beta \epsilon j(j+1)}
\end{equation}

The $2j + 1$ term is due to degeneracy in the energy levels. In the low temperature limit, $k_B T \ll \epsilon$, then $\beta$ is large and we can approximate $Z_1$ using a few terms. 

\begin{equation}
    Z_{rot} = 1 + 3e^{-\beta \epsilon} + 5 e^{-6\beta \epsilon} + 7e^{-12 \beta \epsilon} + \text{ neglect the rest } 
\end{equation}

At $T \rightarrow 0$, then $Z_{rot} \approx 1 \implies U \approx 0$. Excited states are then "frozen out" at low temperatures. 

At high temperatures. Use an integral instead of a sum. Then, 

\begin{equation}
    Z_{rot} = \frac{1}{\beta \epsilon}
\end{equation}

Then,

\begin{equation}
    U_{rot} = -\partial_{\beta}\ln Z_{rot} = k_B T
\end{equation}

This is for $k_B T \gg \epsilon$. This makes sense since we only had 2 degrees of freedom for a rotating spin.

\subsection*{Vibratioanl partition function}

\begin{equation}
    Z_{vib} = \sum_{n}e^{-\beta \epsilon_{vib}}
\end{equation}
Such that 

\begin{equation}
    \epsilon_{vib} = \hbar \omega(n + 1/2), \quad n = 0,1,2,3,\dots
\end{equation}

Then, 

\begin{equation}
    Z_{vib} = \sum_n e^{-\beta\hbar \omega(n + 1/2)} = e^{-\beta \hbar\omega/2}\sum e^{-\beta n\hbar \omega}
\end{equation}

This is a geometric series. Then, the series converges to:

\begin{equation}
    \frac{e^{-\beta \hbar \omega/2}}{1 - e^{-\beta \hbar \omega}} = Z_{vib}
\end{equation}

Then, 

\begin{align}
    \langle \epsilon \rangle &= U = -\partial_{\beta}\ln\left( Z_{vib} \right) \\
    &= \hbar \omega \left( \frac{1}{e^{\beta \hbar \omega} - 1} + \frac{1}{2} \right) \\
    &= \hbar \omega \left( \langle n \rangle + \frac{1}{2} \right)
\end{align}

Then, we can deduce that 

\begin{equation}
    \langle n \rangle = \frac{1}{e^{\beta \hbar \omega} - 1}
\end{equation}

In the low temperature limit, $k_B T \ll \hbar \omega$, then we obtain 

\begin{equation}
    \langle n \rangle \approx e^{-\beta \hbar \omega}
\end{equation}

Then, 

\begin{equation}
    U_{vib} \approx \hbar \omega \left( e^{-\beta \hbar \omega} + \frac{1}{2} \right) \approx \frac{\hbar \omega}{2}
\end{equation}

This is valid in the low temperature limit, and the zero point energy. In the high temperature limit, then $k_B T \gg \hbar \omega$, then 

\begin{equation}
    \langle n \rangle \approx \frac{1}{1 + \beta \hbar \omega - 1} = \frac{k_B T}{\hbar \omega}
\end{equation}

Therefore, 

\begin{equation}
    U_{vib} \approx \hbar \omega \left( \frac{k_B T}{\hbar \omega} + \frac{1}{2} \right) = k_B T + \frac{\hbar \omega}{2} \approx k_B T
\end{equation}

This is valid in the high temperature limit. This is consistent with the equipartition theorem (\textbf{which is only valid at the high temperature limit!}).

Note, 

\begin{equation}
    C_V = \left( \frac{\partial U}{\partial T} \right)_{V}
\end{equation}

For a diatomic gas, then they got a curve that looked liked a positive step function with steps at $3/2$, $5/2$, etc. 

Then, 

\begin{equation}
    C_V^{trans} = \frac{3}{2}N k_B
\end{equation}
\begin{equation}
    C_V^{rot} = Nk_B, \quad k_B T \gg \epsilon
\end{equation}
\begin{equation}
    C_V^{vib} = Nk_B, \quad k_B T \gg \epsilon
\end{equation}

Then, $C_V$ is the sum of all the above. At lower temperatures, only $C_V^{trans}$ mattered. At high temperatures, $C_V$ is the sum of all the above. In between, you still have some rotational and translational, so $5/2$. 

Lecture 6 has the full derivation of the theorem. 
\end{document}